%% Преамбула TeX-файла

% 1. Стиль и язык
\documentclass[utf8x, 14pt]{G7-32} % Стиль (по умолчанию будет 14pt)

% Остальные стандартные настройки убраны в preamble.inc.tex.
\sloppy

% --------------------------------------------------------- %

\usepackage{wrapfig}
\usepackage{float}
\usepackage{microtype}

\usepackage{setspace}

\usepackage{lipsum}

\newcommand{\code}[1]{\texttt{#1}}

\makeatletter
\newcommand{\vhrulefill}[1]
{
	\leavevmode\leaders\hrule\@height#1\hfill \kern\z@
}
\makeatother

% --------------------------------------------------------- %

% Настройки стиля ГОСТ 7-32
% Для начала определяем, хотим мы или нет, чтобы рисунки и таблицы нумеровались в пределах раздела, или нам нужна сквозная нумерация.
\EqInChapter % формулы будут нумероваться в пределах раздела
\TableInChapter % таблицы будут нумероваться в пределах раздела
\PicInChapter % рисунки будут нумероваться в пределах раздела

% Добавляем гипертекстовое оглавление в PDF
\usepackage[
bookmarks=true, colorlinks=true, unicode=true,
urlcolor=black,linkcolor=black, anchorcolor=black,
citecolor=black, menucolor=black, filecolor=black,
]{hyperref}

\AfterHyperrefFix

\usepackage{microtype}% полезный пакет для микротипографии, увы под xelatex мало чего умеет, но под pdflatex хорошо улучшает читаемость

% Тире могут быть невидимы в Adobe Reader
\ifInvisibleDashes
\MakeDashesBold
\fi

\usepackage{graphicx}   % Пакет для включения рисунков

% С такими оно полями оно работает по-умолчанию:
% \RequirePackage[left=20mm,right=10mm,top=20mm,bottom=20mm,headsep=0pt,includefoot]{geometry}
% Если вас тошнит от поля в 10мм --- увеличивайте до 20-ти, ну и про переплёт не забывайте:
\geometry{right=20mm}
\geometry{left=30mm}
\geometry{bottom=20mm}
\geometry{ignorefoot}% считать от нижней границы текста


% Пакет Tikz
\usepackage{tikz}
\usetikzlibrary{arrows,positioning,shadows}

% Произвольная нумерация списков.
\usepackage{enumerate}

% ячейки в несколько строчек
\usepackage{multirow}

% itemize внутри tabular
\usepackage{paralist,array}

%\setlength{\parskip}{1ex plus0.5ex minus0.5ex} % разрыв между абзацами
\setlength{\parskip}{1ex} % разрыв между абзацами
\usepackage{blindtext}

% Центрирование подписей к плавающим окружениям
%\usepackage[justification=centering]{caption}

\usepackage{newfloat}
\DeclareFloatingEnvironment[
placement={!ht},
name=Equation
]{eqndescNoIndent}
\edef\fixEqndesc{\noexpand\setlength{\noexpand\parindent}{\the\parindent}\noexpand\setlength{\noexpand\parskip}{\the\parskip}}
\newenvironment{eqndesc}[1][!ht]{%
    \begin{eqndescNoIndent}[#1]%
\fixEqndesc%
}
{\end{eqndescNoIndent}}

% ---------------------------------------------------------------- %

\usepackage{listings}
\usepackage{listingsutf8}
\lstset{
	basicstyle=\small,
	keywordstyle=\color{blue},
	stringstyle=\color{red},
	commentstyle=\color{gray},
	numbers=left,
	stepnumber=1,
	numberstyle=\tiny,
	numbersep=5pt,
	language=C++,
	tabsize=2,
	frame=single,
	breaklines=true,
	breakatwhitespace=true
}

\newcommand{\imgwc}[4]
{
	\begin{figure}[#1]
		\center{\includegraphics[width=#2]{inc/img/#3}}
		\caption{#4}
		\label{img:#3}
	\end{figure}
}
\newcommand{\imghc}[4]
{
	\begin{figure}[#1]
		\center{\includegraphics[height=#2]{inc/img/#3}}
		\caption{#4}
		\label{img:#3}
	\end{figure}
}
\newcommand{\imgsc}[4]
{
	\begin{figure}[#1]
		\center{\includegraphics[scale=#2]{inc/img/#3}}
		\caption{#4}
		\label{img:#3}
	\end{figure}
}

\titleformat{\section}[block]{\normalfont\bfseries}{\thesection}{14pt}{}
\titleformat{\subsection}[block]{\normalfont\bfseries}{\thesubsection}{14pt}{}
\titleformat{\subsubsection}[block]{\normalfont\bfseries}{\thesubsubsection}{14pt}{}

\titleformat{\chapter}[block]
	{\normalfont\fontsize{14}{14}\bfseries}
	{\thechapter}
	{1em}{}
\titlespacing{\chapter}{\parindent}{0em}{1em}

\titleformat{\section}[block]
	{\normalfont\fontsize{14}{14}\bfseries}
	{\thesection}
	{1em}{}
\titlespacing{\section}{\parindent}{1.5em}{1em}

\titleformat{\subsection}[block]
	{\normalfont\fontsize{14}{14}\bfseries}
	{\thesubsection}
	{1em}{}
\titlespacing{\subsection}{\parindent}{1.5em}{1.5em}

\newcommand{\UnnamedStructuredChapter}[1]
{
	\titleformat{\chapter}[block]
  	{\centering\normalfont\large\bfseries}
  	{\centering\thechapter}
  	{14pt}{\centering\large}
		\chapter*{\MakeUppercase{#1}}
	\titleformat{\chapter}[block]
  	{\normalfont\large\bfseries}
  	{\thechapter}
  	{14pt}{\large}
}

\newcommand{\StructuredChapter}[1]
{
	\titleformat{\chapter}[block]
  	{\centering\normalfont\large\bfseries}
  	{\centering\thechapter}
  	{14pt}{\centering\large}
		\chapter*{\MakeUppercase{#1}}
	\addcontentsline{toc}{chapter}{#1}
	\titleformat{\chapter}[block]
  	{\normalfont\large\bfseries}
  	{\thechapter}
  	{14pt}{\large}
}

\usepackage{caption}
\DeclareCaptionLabelSeparator{bar}{\space---\space}
\captionsetup{labelsep=bar}
\addto\captionsrussian
{
	\def\figurename{Рисунок}
	\def\contentsname{\normalsize СОДЕРЖАНИЕ}
}

\renewcommand{\theenumi}{\arabic{enumi})}
\renewcommand\labelenumi{\theenumi}

\newcommand\labelitemi{\small$\bullet$}
\newcommand\labelitemii{\small$\bullet$}
\newcommand\labelitemiii{\small$\bullet$}

\newcommand{\specialcell}[2][c]
{
  \begin{tabular}[#1]{@{}c@{}}#2\end{tabular}
}

\usepackage{multicol}

\lstset{
  literate={а}{{\selectfont\char224}}1
           {б}{{\selectfont\char225}}1
           {в}{{\selectfont\char226}}1
           {г}{{\selectfont\char227}}1
           {д}{{\selectfont\char228}}1
           {е}{{\selectfont\char229}}1
           {ё}{{\"e}}1
           {ж}{{\selectfont\char230}}1
           {з}{{\selectfont\char231}}1
           {и}{{\selectfont\char232}}1
           {й}{{\selectfont\char233}}1
           {к}{{\selectfont\char234}}1
           {л}{{\selectfont\char235}}1
           {м}{{\selectfont\char236}}1
           {н}{{\selectfont\char237}}1
           {о}{{\selectfont\char238}}1
           {п}{{\selectfont\char239}}1
           {р}{{\selectfont\char240}}1
           {с}{{\selectfont\char241}}1
           {т}{{\selectfont\char242}}1
           {у}{{\selectfont\char243}}1
           {ф}{{\selectfont\char244}}1
           {х}{{\selectfont\char245}}1
           {ц}{{\selectfont\char246}}1
           {ч}{{\selectfont\char247}}1
           {ш}{{\selectfont\char248}}1
           {щ}{{\selectfont\char249}}1
           {ъ}{{\selectfont\char250}}1
           {ы}{{\selectfont\char251}}1
           {ь}{{\selectfont\char252}}1
           {э}{{\selectfont\char253}}1
           {ю}{{\selectfont\char254}}1
           {я}{{\selectfont\char255}}1
           {А}{{\selectfont\char192}}1
           {Б}{{\selectfont\char193}}1
           {В}{{\selectfont\char194}}1
           {Г}{{\selectfont\char195}}1
           {Д}{{\selectfont\char196}}1
           {Е}{{\selectfont\char197}}1
           {Ё}{{\"E}}1
           {Ж}{{\selectfont\char198}}1
           {З}{{\selectfont\char199}}1
           {И}{{\selectfont\char200}}1
           {Й}{{\selectfont\char201}}1
           {К}{{\selectfont\char202}}1
           {Л}{{\selectfont\char203}}1
           {М}{{\selectfont\char204}}1
           {Н}{{\selectfont\char205}}1
           {О}{{\selectfont\char206}}1
           {П}{{\selectfont\char207}}1
           {Р}{{\selectfont\char208}}1
           {С}{{\selectfont\char209}}1
           {Т}{{\selectfont\char210}}1
           {У}{{\selectfont\char211}}1
           {Ф}{{\selectfont\char212}}1
           {Х}{{\selectfont\char213}}1
           {Ц}{{\selectfont\char214}}1
           {Ч}{{\selectfont\char215}}1
           {Ш}{{\selectfont\char216}}1
           {Щ}{{\selectfont\char217}}1
           {Ъ}{{\selectfont\char218}}1
           {Ы}{{\selectfont\char219}}1
           {Ь}{{\selectfont\char220}}1
           {Э}{{\selectfont\char221}}1
           {Ю}{{\selectfont\char222}}1
           {Я}{{\selectfont\char223}}1
}

\usepackage{pdflscape}
\usepackage{rotating}
\usepackage[absolute]{textpos}
\usepackage{fancyhdr}

\usepackage{everypage}
\usepackage{lipsum}


\fancypagestyle{lscape}{% 
\fancyhf{} % clear all header and footer fields 
\fancyfoot[LE]{%
\begin{textblock}{20}(1,5){\rotatebox{90}{\leftmark}}\end{textblock}
\begin{textblock}{1}(13,10.5){\rotatebox{90}{\thepage}}\end{textblock}}
\fancyfoot[LO] {%
\begin{textblock}{1}(13,10.5){\rotatebox{90}{\thepage}}\end{textblock}
\begin{textblock}{20}(1,13.25){\rotatebox{90}{\rightmark}}\end{textblock}}
\renewcommand{\headrulewidth}{0pt} 
\renewcommand{\footrulewidth}{0pt}}

\setlength{\TPHorizModule}{1cm}
\setlength{\TPVertModule}{1cm}

\usepackage{adjustbox}

\newcommand{\Lpagenumber}{\ifdim\textwidth=\linewidth\else\bgroup
  \dimendef\margin=0 %use \margin instead of \dimen0
  \ifodd\value{page}\margin=\oddsidemargin
  \else\margin=\evensidemargin
  \fi
  \raisebox{\dimexpr -\topmargin-\headheight-\headsep-0.5\linewidth}[0pt][0pt]{%
    \rlap{\hspace{\dimexpr \margin+\textheight+\footskip}%
    \llap{\rotatebox{90}{\thepage}}}}%
\egroup\fi}
\AddEverypageHook{\Lpagenumber}%

% Настройки листингов.
\ifPDFTeX
% 8 Листинги

\usepackage{listings}

% Значения по умолчанию
\lstset{
  basicstyle= \footnotesize,
  breakatwhitespace=true,% разрыв строк только на whitespacce
  breaklines=true,       % переносить длинные строки
%   captionpos=b,          % подписи снизу -- вроде не надо
  inputencoding=koi8-r,
  numbers=left,          % нумерация слева
  numberstyle=\footnotesize,
  showspaces=false,      % показывать пробелы подчеркиваниями -- идиотизм 70-х годов
  showstringspaces=false,
  showtabs=false,        % и табы тоже
  stepnumber=1,
  tabsize=4,              % кому нужны табы по 8 символов?
  frame=single
}

% Стиль для псевдокода: строчки обычно короткие, поэтому размер шрифта побольше
\lstdefinestyle{pseudocode}{
  basicstyle=\small,
  keywordstyle=\color{black}\bfseries\underbar,
  language=Pseudocode,
  numberstyle=\footnotesize,
  commentstyle=\footnotesize\it
}

% Стиль для обычного кода: маленький шрифт
\lstdefinestyle{realcode}{
  basicstyle=\scriptsize,
  numberstyle=\footnotesize
}

% Стиль для коротких кусков обычного кода: средний шрифт
\lstdefinestyle{simplecode}{
  basicstyle=\footnotesize,
  numberstyle=\footnotesize
}

% Стиль для BNF
\lstdefinestyle{grammar}{
  basicstyle=\footnotesize,
  numberstyle=\footnotesize,
  stringstyle=\bfseries\ttfamily,
  language=BNF
}

% Определим свой язык для написания псевдокодов на основе Python
\lstdefinelanguage[]{Pseudocode}[]{Python}{
  morekeywords={each,empty,wait,do},% ключевые слова добавлять сюда
  morecomment=[s]{\{}{\}},% комменты {а-ля Pascal} смотрятся нагляднее
  literate=% а сюда добавлять операторы, которые хотите отображать как мат. символы
    {->}{\ensuremath{$\rightarrow$}~}2%
    {<-}{\ensuremath{$\leftarrow$}~}2%
    {:=}{\ensuremath{$\leftarrow$}~}2%
    {<--}{\ensuremath{$\Longleftarrow$}~}2%
}[keywords,comments]

% Свой язык для задания грамматик в BNF
\lstdefinelanguage[]{BNF}[]{}{
  morekeywords={},
  morecomment=[s]{@}{@},
  morestring=[b]",%
  literate=%
    {->}{\ensuremath{$\rightarrow$}~}2%
    {*}{\ensuremath{$^*$}~}2%
    {+}{\ensuremath{$^+$}~}2%
    {|}{\ensuremath{$|$}~}2%
}[keywords,comments,strings]

% Подписи к листингам на русском языке.
\renewcommand\lstlistingname{Листинг}
\renewcommand\lstlistlistingname{Листинги}

\else
\usepackage{local-minted}
\fi

% Полезные макросы листингов.
% Любимые команды
\newcommand{\Code}[1]{\textbf{#1}}


\begin{document}

\frontmatter % выключает нумерацию ВСЕГО; здесь начинаются ненумерованные главы: реферат, введение, глоссарий, сокращения и прочее.

% Стиль титульного листа и заголовки
%\thispagestyle{empty}

\begin{titlepage}
	\centering
	
	\begin{wrapfigure}[7]{l}{0.14\linewidth}
		\vspace{5mm}
		\hspace{-5.8mm}
		\includegraphics[width=0.93\linewidth]{inc/img/bmstu-logo}
	\end{wrapfigure}
	{\singlespacing \footnotesize \bfseries Министерство науки и высшего образования Российской Федерации\\Федеральное государственное бюджетное образовательное учреждение\\высшего образования\\<<Московский государственный технический университет\\имени Н.~Э.~Баумана\\ (национальный исследовательский университет)>>\\(МГТУ им. Н.~Э.~Баумана)\\}

	\vspace{-2.2mm}
	\vhrulefill{0.9mm}\\
	\vspace{-7mm}
	\vhrulefill{0.2mm}\\
	\vspace{2mm}
	
	{\doublespacing \small \raggedright ФАКУЛЬТЕТ \hspace{25mm} «Информатика и системы управления»\\
	КАФЕДРА \hspace{5mm} «Программное обеспечение ЭВМ и информационные технологии»\\}

	\vspace{30mm}
	
	\textbf{ОТЧЕТ}\\
	По лабораторной работе №2\\
	По курсу: Анализ алгоритмов\\
	Тема: Алгоритмы умножения матриц\\
	%Вариант 0\\

	\vspace{60mm}

	\hspace{70mm} Студент:       \hfill Лисневский~А.~В.\\
	\hspace{70mm} Группа:        \hfill ИУ7-52Б\\
	\hspace{70mm} Преподаватель: \hfill Волкова~Л.~Л.\\
	\hspace{70mm} Оценка:        \hfill \hrulefill\\

	\vfill
	
	\begin{center}
		Москва, \the\year
	\end{center}
\end{titlepage}

\setcounter{page}{1}

\setcounter{page}{2}

\tableofcontents

\StructuredChapter{Введение}

	Целью данной работы является изучение архитектуры гетерогенных вычислительных систем и технологии разработки ускорителей вычислений на базе ПЛИС фирмы Xilinx.

Для достижения данной цели необходимо выполнить следующие задачи:
\begin{itemize}
	\item изучить основные сведения о платформе Xilinx Alveo U200;
	\item разработать RTL описание ускорителя вычислений по индивидуальному варианту;
	\item выполнить генерацию ядра ускорителя;
	\item выполнить синтез и сборку бинарного модуля ускорителя;
	\item разработать и отладить тестирующее программное обеспечение на серверной хост-платформе;
	\item провести тесты работы ускорителя вычислений.
\end{itemize}


\mainmatter % это включает нумерацию глав и секций в документе ниже

\setcounter{chapter}{1}

\section*{Задание 1}
\addcontentsline{toc}{section}{Задание 1}

\textbf{Задание}: Ознакомиться с возможностями программы PCLAB в Разделе 2 методических указаний. Запустить программу PCLAB 1.0. Изучить идентификационную информацию на вкладке «Идентификация процессора». 

Описание PCLAB: Программа PCLAB предназначена для исследования производительности x86 совместимых ЭВМ c IA32 архитектурой, работающих под управлением операционной системы Windows (версий 95 и старше). Исследование организации ЭВМ заключается в проведении ряда экспериментов, направленных на построение зависимостей времени обработки критических участков кода от изменяемых параметров. 

Набор реализуемых программой экспериментов позволяет исследовать особенности построения современных подсистем памяти ЭВМ и процессорных устройств, выявить конструктивные параметры конкретных моделей ЭВМ. Процесс сбора и анализа экспериментальных данных в PCLAB основан на процедуре профилировки критического кода, т.е. в измерении времени его обработки центральным процессорным устройством.

\section*{Задание 2}
\addcontentsline{toc}{section}{Задание 2}

\textbf{Задание}: На основании идентификационной информации о микропроцессоре ЭВМ, используемой при проведении лабораторной работы, определить следующие параметры: размер линейки кэш-памяти верхнего уровня и объем физической памяти. Результаты занести в отчет.

Размер линейки кэша: 64 байт;
Объем физической памяти: 4 Гб.

\section*{Задание 3}
\addcontentsline{toc}{section}{Задание 3}

\textbf{Задание}: Ознакомиться с описанием эксперимента «Исследование расслоения динамической памяти» на вкладке «Описание эксперимента». Провести эксперимент. По результатам эксперимента определить: количество банков динамической памяти; размер одной страницы динамической памяти; количество страниц в динамической памяти. Сделать выводы о использованном способе наращивания динамической памяти. Результаты занести в отчет.

\textbf{Цель эксперимента}: определение способа трансляции физического адреса, используемого при обращении к динамической памяти.

\textbf{Исходные данные}:
\begin{itemize}
	\item Размер линейки кэша: 64 байт;
	\item Объем физической памяти: 4 Гбайт.
\end{itemize}

В таблице \ref{tab_3} приведены настраиваемые параметры.
\begin{table}[H]
	\begin{center}
		\caption{Настраиваемые параметры}
		\label{tab_3}
		\begin{tabular}{|c|c|c|c|}
		\hline
		Эксперимент & № & Значение & Описание 	\\
		\hline
		\hline
		\multirow{3}*{1} & 1 & 32 Кбайт & \specialcell{Максимальное расстояния \\ между читаемыми блоками} 	\\
		\cline{2-4} & 2 & 128 байт &\specialcell{ Шаг увеличения расстояния между \\ читаемыми 4-х байтовыми ячейками}	\\
		\cline{2-4} & 3 & 1 Мбайт & Размер массива		\\
		\hline
		\multirow{3}*{2} & 1 & 32 Кбайт & \specialcell{Максимальное расстояния \\ между читаемыми блоками} 	\\
		\cline{2-4} & 2 & 64 байт &\specialcell{ Шаг увеличения расстояния между \\ читаемыми 4-х байтовыми ячейками}	\\
		\cline{2-4} & 3 & 1 Мбайт & Размер массива		\\
		\hline
		\end{tabular}
	\end{center}
\end{table}

Ниже, на рисунках \ref{img:task_01} - \ref{img:task_11} приведены зависимости времени обращения к памяти от расстояния между читаемыми блоками данных.

\imgsc{H}{0.4}{task_01}{Результаты исследования расслоения динамической памяти (часть 1)}
\imgsc{H}{0.4}{task_21}{Результаты исследования расслоения динамической памяти (часть 2)}
\imgsc{H}{0.4}{task_11}{Результаты исследования расслоения динамической памяти (часть 3)}

Ниже, на рисунке \ref{img:bank_memory} приведена блочно-циклическая схема расслоения памяти, помогающая более ясно представить суть эксперимента.

\imgsc{H}{0.85}{bank_memory}{Блочно-циклическая схема расслоения памяти}

\textbf{Результаты эксперимента}:
\begin{itemize}
	\item количество банков динамической памяти: $B = T1/P = 1024/128 = 8$;
	\item размер одной страницы динамической памяти: $PS = T2/B = 4096/8 = 512\ byte$;
	\item количество страниц в динамической памяти: $C = V/(PS*B*P) = \frac{2^{32}}{2^{19}} = 8192\ byte$.
\end{itemize}

Примечание:
\begin{itemize}
	\item B --- количество банок памяти;
	\item P --- объем данных, являющийся минимальной порцией обмена кэш-памяти верхнего уровня с оперативной памятью;
	\item PS --- размер страницы DRAM памяти;
	\item T1 --- размер блока в одной банке памяти;
	\item T2 --- размер страницы одной банки памяти;
	\item V --- объем физического пространства ОП;
	\item C --- количество страниц в ОП.
\end{itemize}

\textbf{Выводы}:
\begin{enumerate}
	\item Память расслоена, доступ к ней с разным временем доступа, в зависимости от размера запрашиваемого блока.
	\item Экспериметально были получены значения количества банков ОП, размер страница банка, размер страницы DRAM.
	\item Данные следует размещать так, чтобы они укладывали в одну страницу.
	\item Данные следует выравнивать по размеру линейки кэша.
	\item Данные следует обрабатывать так, чтобы минимизировать количество последовательных обращений в одну банку ОП.
\end{enumerate} 

\section*{Задание 4}
\addcontentsline{toc}{section}{Задание 4}

\textbf{Задание}: Ознакомиться с описанием эксперимента «Сравнение эффективности ссылочных и векторных структур данных». Провести эксперимент. По результатам эксперимента определить: отношение времени работы алгоритма, использующего зависимые данные, ко времени обработки аналогичного алгоритма обработки независимых данных. Сделать выводы об эффективности ссылочных и векторных структур данных и способах ее повышения. Результаты занести в отчет.

\textbf{Цель эксперимента}: оценить влияние зависимости команд по данным на эффективность вычислений.

В таблице \ref{tab_4} приведены настраиваемые параметры.
\begin{table}[H]
	\begin{center}
		\caption{Настраиваемые параметры}
		\label{tab_4}
		\begin{tabular}{|c|c|c|c|}
		\hline
		№ & Значение & Описание 	\\
		\hline
		\hline
		1 & 1 Мбайт & \specialcell{Количество элементов в списке} 	\\
		\hline
		2 & 32 Кбайт & \specialcell{Максимальная фрагментации списка}		\\
		\hline
		3 & 1 Кбайт & \specialcell{Шаг увеличения фрагментации}	\\
		\hline
		\end{tabular}
	\end{center}
\end{table}

Ниже, на рисунке \ref{img:task_02} приведена зависимость времени выполнения поиска минимального значения для массива и односвязного списка.

\imgsc{H}{0.41}{task_02}{Сравнение эффективности ссылочных и векторных структур данных}

Из полученного графика видна проблема сематического разрыва. Следует использовать структуры данных с учётом технологического фактора. Для машины «лучше» использовать массив т.к. она плохо работает со списками.

\textbf{Результаты эксперимента}. Односвязный список обрабатывался в \textbf{19,741378} раз дольше. \\

\textbf{Выводы}:
\begin{enumerate}
	\item Связанные данные следует организовывать так, чтобы при работе программы они были как можно ближе друг к другу расположены в ОП.
	\item Использование структур данных, помогающих ясней представить и быстрее решить задачу, может приводить к значительному снижению производительности системы (Таким образом, в эксперименте наблюдали следствие сематического разрыва)
\end{enumerate} 

\section*{Задание 5}
\addcontentsline{toc}{section}{Задание 5}

\textbf{Задание}: Для ЭВМ, используемой при проведении лабораторной работы определить следующие параметры: степень ассоциативности и размер TLB данных. Ознакомиться с описанием и провести эксперимент «Исследование эффективности программной предвыборки». По результатам эксперимента определить: отношение времени последовательной обработки блока данных ко времени обработки блока с применением предвыборки; время и количество тактов первого обращения к странице данных. Сделать выводы об эффективности предвыборки и способах ее повышения. Результаты занести в отчет.

\textbf{Цель эксперимента}: выявить способы ускорения вычислений благодаря применению предвыборки данных.

\textbf{Исходные данные}: 
\begin{itemize}
	\item степень ассоциативности TLB данных: 4 ячейки;
	\item размер TLB данных: 128 групп.
\end{itemize}

В таблице \ref{tab_5} приведены настраиваемые параметры.
\begin{table}[H]
	\begin{center}
		\caption{Настраиваемые параметры}
		\label{tab_5}
		\begin{tabular}{|c|c|c|c|}
		\hline
		№ & Значение & Описание 	\\
		\hline
		\hline
		1 & 512 байт & \specialcell{Шаг увеличения расстояния \\ между читаемыми данными} \\
		\hline
		2 & 128 Кбайт & \specialcell{Размер массива}		\\
		\hline
		\end{tabular}
	\end{center}
\end{table}

Ниже, на рисунке \ref{img:task_03} приведены зависимости времени обращения к памяти от расстояния между читаемыми блоками данных. Красный график показывает время или количество тактов работы алгоритма без предвыборки. Зеленый график показывает время или количество тактов работы алгоритма с использованием предвыборки.

\imgsc{H}{0.4}{task_03}{Исследование эффективности программной предвыборки}

Время обращения к первому элементу в таблицы в 20 раз больше т.к. мы не знаем где это страница находится в памяти, мы имеем только логический адрес, а нужен физический. Поэтому логично использовать предвыборку. В данном примере ускорение почти в 2 раза. 

\textbf{Результаты эксперимента}:
\begin{itemize}
	\item обработка без загрузки таблицы страниц в TLB производилась в \textbf{1,8456621} раз дольше;
	\item время первого обращения к странице данных: \textbf{15`000} тактов.
\end{itemize}

\textbf{Вывод}: Для исключения задержек, связанных с получением физического адреса начала страницы, имеет смысл предварительно загрузить страницы в TLB перед работой с большими массивами данных.

\section*{Задание 6}
\addcontentsline{toc}{section}{Задание 6}

\textbf{Задание}: Ознакомиться с описанием и провести эксперимент «Исследование способов эффективного чтения оперативной памяти». По результатам эксперимента определить: отношение времени обработки блока памяти неоптимизированной структуры ко времени обработки блока структуры, обеспечивающей эффективную загрузку и параллельную обработку данных. Сделать выводы о способах повышения эффективности чтения оперативной памяти.

\textbf{Цель эксперимента}: исследование возможности ускорения вычислений благодаря
использованию структур данных, оптимизирующих механизм чтения оперативной памяти.

\textbf{Исходные данные}: 
\begin{itemize}
	\item адресное расстояние между банками памяти: 128 байт;
	\item размер буфера чтения:	4 КБайт.
\end{itemize}

В таблице \ref{tab_6} приведены настраиваемые параметры.
\begin{table}[H]
	\begin{center}
		\caption{Настраиваемые параметры}
		\label{tab_6}
		\begin{tabular}{|c|c|c|c|}
		\hline
		№ & Значение & Описание 	\\
		\hline
		\hline
		1 & 1 Мбайт & \specialcell{Размер массива} \\
		\hline
		2 & 128 ед. & \specialcell{Количество потоков данных}		\\
		\hline
		\end{tabular}
	\end{center}
\end{table}

Ниже, на рисунке \ref{img:task_04} приведены зависимости времени чтения данных от количества одновременно обрабатываемых массивов для неоптимизированной структуры (красный график) и структуры, обеспечивающей эффективную загрузку и параллельную обработку данных (зеленый график).

\imgsc{H}{0.4}{task_04}{Исследование способов эффективного чтения оперативной памяти}

\textbf{Результаты эксперимента}:
\begin{itemize}
	\item отношение времени обработки блока памяти неоптимизированной структуры ко времени обработки блока структуры, обеспечивающей эффективную загрузку и параллельную обработку данных: \textbf{1.65}.
\end{itemize}

\textbf{Вывод}: 
\begin{itemize}	
	\item упорядочив данные определённым образом, можно ускорить приложение (мы группируем данные, которые используем вместе.);
	\item следует переупорядочивать данные, выравнивая их по размеру кэш-линии, тем самым исключая несвоевременную передачу данных;
	\item следует размещать данные как можно ближе друг к другу: стараться, по возможности, не обращаться к диспетчеру кучи за памятью, помнить и использовать на приктике особенности  выравнивания данных (например в С-структурах);
\end{itemize}

\section*{Задание 7}
\addcontentsline{toc}{section}{Задание 7}

\textbf{Задание}: Для ЭВМ, используемой при проведении лабораторной работы определить следующие параметры: размер банка кэш-памяти данных первого и второго уровня, степень ассоциативности кэш-памяти первого и второго уровня, размер линейки кэш-памяти первого и второго уровня. Ознакомиться с описанием и провести эксперимент «Исследование конфликтов в кэш-памяти». По результатам эксперимента определить: отношение времени обработки массива с конфликтами в кэш-памяти ко времени обработки массива без конфликтов. Сделать выводы о способах устранения конфликтов в кэш-памяти. 

\textbf{Цель эксперимента}: исследование влияния конфликтов кэш-памяти на эффективность
вычислений.

\textbf{Исходные данные}:
\begin{itemize}
	\item размер банка кэш-памяти данных первого и второго уровня: 32 КБайт;
	\item степень ассоциативности кэш-памяти первого и второго уровня: 8 ячеек;
	\item размер линейки кэш памяти первого и второго уровня:	64 байт.
\end{itemize}

В таблице \ref{tab_7} приведены настраиваемые параметры.
\begin{table}[H]
	\begin{center}
		\caption{Настраиваемые параметры}
		\label{tab_7}
		\begin{tabular}{|c|c|c|c|}
		\hline
		№ & Значение & Описание 	\\
		\hline
		\hline
		1 & 128 Кбайт & \specialcell{Размер банка кэш-памяти} \\
		\hline
		2 & 128 байт & \specialcell{Размер линейки кэш-памяти}		\\
		\hline
		3 & 32 ед. & \specialcell{Количество читаемых линеек}		\\
		\hline
		\end{tabular}
	\end{center}
\end{table}

Ниже, на рисунке \ref{img:task_05} приведены зависимости времени обращения к памяти от расстояния между читаемыми блоками данных.

\imgsc{H}{0.4}{task_05}{Исследование конфликтов в кэш-памяти}

\textbf{Результаты эксперимента}: отношение времени обработки массива с конфликтами в кэш-памяти ко времени обработки массива без конфликтов = \textbf{5,9547263}.

Ниже, на рисунке \ref{img:cache} приведена кэш-память с четырехканальным частично-ассоциативным отображением, помогающая понять суть эксперимента.

\imgsc{H}{0.75}{cache}{Кэш-память с четырехканальным частично-ассоциативным отображением}

\textbf{Вывод}:
\begin{itemize}
	\item кэш память ускоряет процессор примерно в 6 раз;
	\item данные следует выравнивать по адресам, кратным размеру линейки кэша;
	\item данные следует обрабатывать так, чтобы уменьшить количество последовательных обращений к блокам памяти, соответствующих одному набору (модулю) (рисунок \ref{img:cache}).
\end{itemize}

\section*{Задание 8}
\addcontentsline{toc}{section}{Задание 8}

\textbf{Задание}: Ознакомиться с описанием и провести эксперимент «Исследование алгоритмов сортировки». По результатам эксперимента определить: отношение времени сортировки массивов алгоритмом QuickSort ко времени сортировки алгоритмом Counting Radix, а также ко времени сортировки Counting-Radix алгоритмом, оптимизированным под 8-процессорную вычислительную систему. Сделать выводы о наиболее эффективном алгоритме сортировки.

\textbf{Цель эксперимента}: исследование способов эффективного использования памяти и
выявление наиболее эффективных алгоритмов сортировки, применимых в вычислительных
системах.

В таблице \ref{tab_8} приведены настраиваемые параметры.
\begin{table}[H]
	\begin{center}
		\caption{Настраиваемые параметры}
		\label{tab_8}
		\begin{tabular}{|c|c|c|c|}
		\hline
		№ & Значение & Описание 	\\
		\hline
		\hline
		1 & 1 Мбайт & \specialcell{Количество 64-х разрядных \\ элементов массивов} \\
		\hline
		2 & 8 Кбайт & \specialcell{Шаг увеличения размера массива}		\\
		\hline
		\end{tabular}
	\end{center}
\end{table}

Ниже, на рисунке \ref{img:task_06} приведены зависимости времени сортировки (алгоритмы Quick Sort, Radix-Counting Sort, Оптимизированный Radix-Counting Sort) от размера исходного массива.

\imgsc{H}{0.4}{task_06}{Исследование алгоритмов сортировки}

\textbf{Примечание}: Фиолетовый график показывает время или количество тактов работы алгоритма QuickSort. Красный график показывает время или количество тактовработы неоптимизированного алгоритма Radix-Counting. Зеленый график показывает время или количество тактов работы оптимизированного под 8-процессорную вычислительную систему алгоритма Radix-Counting.

\textbf{Результаты эксперимента}:
\begin{itemize}
	\item отношение времени сортировки массива алгоритмом QuickSort ко времени сортировки алгоритмом Radix-Counting Sort: 	\textbf{1.8379134}; 
	\item отношение времени сортировки массива алгоритмом QuickSort ко времени сортировки
Radix-Counting Sort, оптимизированной под 8-процессорную вычислительную систему: \textbf{2.0982673}.
\end{itemize}

\textbf{Выводы}: 
\begin{itemize}
	\item существует алгоритм поразрядной сортировки с сложности меньше чем линейной вычислительной сложности O(n/log(n));
	\item следует выбирать алгоритмы на основе типа входных данных, которые позволяют решить задачу наиболее эффективно.
\end{itemize}
\chapter{Задание 1}
\label{cha:design}

\begin{lstlisting}[label=code:source_1, caption=Листинг исходной программы]
        .section .text
        .globl _start;
        len = 8 
        enroll = 4
		elem_sz = 4
_start:
        la x1, _x
        addi x20, x1, elem_sz*len #Адрес последнего элемента
lp:
        lw x2, 0(x1)
        add x31, x31, x2 #!
        lw x3, 4(x1)
        add x31, x31, x3
        lw x4, 8(x1)
        add x31, x31, x4
        lw x5, 12(x1)
        add x31, x31, x5
        addi x1, x1, elem_sz*enroll
        bne x1, x20, lp
        addi x31, x31, 1
lp2: j lp2

        .section .data
_x:     .4byte 0x1
        .4byte 0x2
        .4byte 0x3
        .4byte 0x4
        .4byte 0x5
        .4byte 0x6
        .4byte 0x7
        .4byte 0x8
\end{lstlisting}

\begin{lstlisting}[label=code:source_21, caption=Деассемблированный листинг исходной программы]
Disassembly of section .text:

80000000 <_start>:
80000000:       00000097                auipc   x1,0x0
80000004:       03c08093                addi    x1,x1,60 # 8000003c <_x>
80000008:       02008a13                addi    x20,x1,32
\end{lstlisting}
\begin{lstlisting}[label=code:source_22, caption=Продолжение листинга \ref{code:source_21}]
8000000c <lp>:
8000000c:       0000a103                lw      x2,0(x1)
80000010:       002f8fb3                add     x31,x31,x2
80000014:       0040a183                lw      x3,4(x1)
80000018:       003f8fb3                add     x31,x31,x3
8000001c:       0080a203                lw      x4,8(x1)
80000020:       004f8fb3                add     x31,x31,x4
80000024:       00c0a283                lw      x5,12(x1)
80000028:       005f8fb3                add     x31,x31,x5
8000002c:       01008093                addi    x1,x1,16
80000030:       fd409ee3                bne     x1,x20,8000000c <lp>
80000034:       001f8f93                addi    x31,x31,1

80000038 <lp2>:
80000038:       0000006f                jal     x0,80000038 <lp2>
\end{lstlisting}

Ниже, в листинге \ref{code:source_31} приведен псевдокод на языке C эквивалентной программы.

\begin{lstlisting}[label=code:source_31, caption=Псевдокод на языке C эквивалентной программы]
#define len 8
#define enroll 4
#define elem_sz 4

int _x[] = {1,2,3,4,5,6,7,8};
int x2, x3, x4, x5, x31;

void _start()
{
	int *x1 = _x;
    int x20 = len;

    do
    {
      x2 = x1[0];
      x31 += x2;
      x3 = x1[1];
      x31 += x3;
      x4 = x1[2];
      x31 += x4;
      x5 = x1[3];
      x31 += x5;
      x1 += enroll;
\end{lstlisting}
\begin{lstlisting}[label=code:source_32, caption=Продолжение листинга \ref{code:source_31}]
    } while(x1 != x20);
    x31++;

    while(1){}
}
\end{lstlisting}

В результате выполнения данного участка кода в регистр x31 будет находится инкрементированная  сумма элементов массива, то есть значение 0x25.

\chapter{Задания 2, 3, 4}
\label{cha:impl}

Мой вариант 17, следовательно в соответствии с заданием $2$ выполняем поиск такта, в котором выполняется выборка команды с адрессом $0x80000020$, номер итерации: 2-я. на рисунке \ref{img:FetchCompleteMy} приведена временная диаграмма, поясняющая этапы выборки и диспетчеризации.

\imgsc{H}{0.6}{FetchCompleteMy}{Временнуя диаграмма выполнения стадий выборки и диспетчеризации команды}

В соответствии с заданием 3 выполняем поиск такта, в котором выполняется декодирование и планирование команды с адрессом $0x8000002$с, номер итерации: 2-я. На рисунке \ref{img:DecodeAndIssue} приведена временная диаграмма, поясняющая этап декодирования и планирования.

\imgsc{H}{0.6}{DecodeAndIssue}{Временнуя диаграмма выполнения стадии декодирования и планирования на выполнение команды}

В соответствии с заданием 4 выполняем поиск такта, в котором выполняется исполнение команды с адрессом $0x80000018$, номер итерации: 2-я. На рисунке \ref{img:Execute} приведена временная диаграмма, поясняющая этап выполнения команды.

\imgsc{H}{0.6}{Execute}{Диаграмма, соответствующая этапу выполнения}

\chapter{Задание 5}
На рисунках \ref{img:fetchANDdispatch5}-\ref{img:execute5} изображены временные диаграммы сигналов, соответствующих всем стадиям выполнения команды, обозначенной в тексте программы символом \#!.

\imgsc{H}{0.6}{fetchANDdispatch5}{Диаграмма, соответствующая этапам выборки и диспетчеризации команды, обозначенной в тексте программы символом \#!}

\imgsc{H}{0.6}{decodeANDissue5}{Диаграмма, соответствующая этапам декодирования и планирования команды, обозначенной в тексте программы символом \#!}

\newpage
\imgsc{H}{0.6}{execute5}{Диаграмма, соответствующая этапу выполнения команды, обозначенной в тексте программы символом \#!}

Результат выполнения программы занесен в регистр x31. Как и предполагалось, в нем хранится значение 0x31 на момент окончания работы программы (рисунок \ref{img:rez}).

\imgsc{H}{0.6}{rez}{Результат выполнения программы}
\chapter{Индивидуальное задание}
\label{cha:research}

\section{Исходная программа}

Вернемся к изучению программы, выданной по индивидуальному заданию. Ее содержание отличиется тем, что вместо уменьшения счетчика цикла, увеличивается адрес блока значений для суммирования. 

Более наглядно порядок выполнения стадий конвейера процессора можно посмотреть с использованием диаграммы трассы работы программы.

Ниже, на рисунке \ref{img:pipeline_1} приведена трасса работы программы.

\imgsc{H}{0.40}{pipeline_1}{Диаграмма, соответствующая этапам выборки и диспетчеризации}

Примечание --- цифрой $1$ обозначен этап выборки (адрес интересующей команды равен $0x80000018$), данной команде на данном этапе присваивается $pc\_id = 6$; цифрой 2 обозначен этап диспетчеризации; цифрой 3 обозначен этап дешифрования и планирования (декодирование выполняется не сразу, так как сигнал $decode\_advance = 0$); цифрой 4 обозначен этап выполения (этап выполняется не сразу из-за конфликта [сигнал $rs2\_conflict = 1$]); во время такта, на котором остановлен желтый селектор, уже виден результат выполнения ($unit\_wb[0]/new\_request = 1$; $unit\_wb[0].done = 1$; $unit\_wb[0].id = 6$)

Ниже, на рисунке \ref{img:task_051} приведена временная диаграмма, поясняющая этапы обработки команды $add x31, x31, x2 \#!$.

\imgsc{H}{0.50}{task_051}{Диаграмма, поясняющая обработку команды $add x31, x31, x2 \#!$}

\section{Оптимизация программы}

Из рассмотрения выше представленной трассы видно, что выполнение программы до окончания последней команды $addi x1, x1, 16$ заняло 45 тактов (1-45). Из них, в течение 18 тактов(
7, 8, 11, 12, 15, 16, 19, 20, 25, 26, 29, 30, 33, 34, 37, 38, 41, 42) не происходило планирование на выполнение новых команд, то есть, 40\% времени было потрачено не эффективно из-за ошибок предсказания ветвлений и конфликтов на конвейере. Имеется возможность существенно сократить время выполнения путем перестановки команд для ликвидации конфликтов. 

Ниже в листинге \ref{code:source_2} приведена оптимизированная программа. Доступ к памяти выполняется за 3 такта, поэтому можно переставить команды так, чтобы регистры, в которые загружаются элементы массива блоком обращения к памяти, были бы освобождены за один такт до того, как с ними начнет работать АЛУ. Таким образом, можно сделать предположение, что для уменьшения времени работы программы загрузку всех необходимых данных в регистры из физической памяти следует выполнять до выполнения каких-либо арифметических операций.

\begin{lstlisting}[label=code:source_2, caption=Листинг оптимизированной программы]
		.section .text
_start:
        la x1, _x
        addi x20, x1, elem_sz*len #Адрес последнего элемента
lp:
        lw x2, 0(x1)
        lw x3, 4(x1)
        lw x4, 8(x1)
        lw x5, 12(x1)
        add x31, x31, x2
        add x31, x31, x3
        add x31, x31, x4
        add x31, x31, x5
        addi x1, x1, elem_sz*enroll
        bne x1, x20, lp
        addi x31, x31, 1
lp2: j lp2
        .section .data
        ...
\end{lstlisting}

\begin{lstlisting}[label=code:source_2, caption=Деассемблированный листинг оптимизированной  программы]
Disassembly of section .text:

80000000 <_start>:
80000000:       00000097                auipc   x1,0x0
80000004:       03c08093                addi    x1,x1,60 # 8000003c <_x>
80000008:       02008a13                addi    x20,x1,32

8000000c <lp>:
8000000c:       0000a103                lw      x2,0(x1)
80000010:       0040a183                lw      x3,4(x1)
80000014:       0080a203                lw      x4,8(x1)
80000018:       00c0a283                lw      x5,12(x1)
8000001c:       002f8fb3                add     x31,x31,x2
80000020:       003f8fb3                add     x31,x31,x3
80000024:       004f8fb3                add     x31,x31,x4
80000028:       005f8fb3                add     x31,x31,x5
8000002c:       01008093                addi    x1,x1,16
80000030:       fd409ee3                bne     x1,x20,8000000c <lp>
80000034:       001f8f93                addi    x31,x31,1

80000038 <lp2>:
80000038:       0000006f                jal     x0,80000038 <lp2>
\end{lstlisting}

Ниже, на рисунке \ref{img:pipeline_2} приведена трасса работы оптимизированной программы.

\imgsc{H}{0.50}{pipeline_2}{Трасса работы оптимизированной программы}

\section{Вывод}

Выполнение участка кода до последней команды $addi x1, x1, 16$ заняло 29 тактов, что в 1.55 раза меньше времени выполнения неоптимизированной программы. Из них только втечение 2 тактов не происходило планирование на выполнение новых команд (около 7\%). Из этого можно заключить, что предложенная оптимизация оправдана.

\backmatter %% Здесь заканчивается нумерованная часть документа и начинаются ссылки и
            
\newpage

\section{Ответы на контрольные вопросы}

\textbf{1. Назовите причины расслоения оперативной памяти}

Расслоение памяти позволяет повысить пропускную способность оперативной памяти (ОП) за счет компоновки ОП из нескольких банков. При таком построении, процедуры обращений к нескольким банкам памяти можно совместить.

\textbf{2. Как в современных процессорах реализована аппаратная предвыборка?}

Если говорить о суперскалярных процессорах, то суть предвыборки заключается в том, чтобы как можно лучше загрузить вычислительный конвейер: при каждой возможности производится считывание команд из памяти, опережающее ход вычислений. Команды после выборки помещаются в быстродействующем буфере предвыборки. Буфер, например, может быть организован по принципу очереди, тогда команды в него поступают в порядке их выполнения в программе. Высокое быстродействие буфера и наличие в нем значительного количества команд позволяет одновременно загружать все конвейеры процессора.

Также можно привести другой пример аппаратной предвыборки: уменьшение времени получения физического адреса страницы ОП при загрузке его в TLB. Суть предвыборки заключается в загрузке физического адреса не одной страницы, а сразу нескольких (8 или 16) за одно обращение к памяти [это возможно за счет расслоения памяти].

\textbf{3. Какая информация храниться в TLB?}

Если рассматривать машину со страничной организацией виртуальной памяти, то физические адреса начала страниц в физической памяти (фреймов).

\textbf{4. Какой тип ассоциативной памяти используется в кэш-памяти второго уровня современных ЭВМ и почему?}

Частично-ассоциативный, так как сочетает  достоинства  прямого  и  полностью  ассоциативного  способов отображения: экономия памяти тэгов (по сравнению с полностью ассоциативным кэшем) и уменьшение количества попаданий в один и тот же набор линеек кэша (по сравнению с прямым отображением).

\textbf{5. Приведите пример программной предвыборки}

Ниже в листинге \ref{code:code_5} приведена часть программы из 3 эксперимента.

\begin{lstlisting}[label=code:code_5, caption=Пример программной предвыборки]
int *p = get_array(arr_size);	// Get huge array from enternal world
int tmp;
for (int a = 0; a < arr_size; a += SIZEOF_PAGE)
	tmp = *(int *) (int(p) + a);

for (int a = 0; a < arr_size; a += 1)
{
	int elem = *(int *) (int(p) + a);
	...							// Array processing
}
\end{lstlisting} %% заключение


%\addcontentsline{toc}{chapter}{Список использованных источников}
\UnnamedStructuredChapter{Список использованных источников}
\bibliographystyle{utf8gost705u}
\begin{thebibliography}{1}
	\bibitem{3}
	Scharstein D., Szeliski R. A taxonomy and evaluation of dense
	two-frame stereo correspondence algorithms // Int. Journal of
	Computer Vision 47. April-June 2002. С. 7–42.
	
	\bibitem{4}
	Bobick A. F., Intille S. S. Large occlusion stereo // Int. Journal
	of Computer Vision. 33(3). 1999. С. 181–200.
	
	\bibitem{5}
	Ohta Y., Kanade T. Stereo by intra- and inter- scanline search
	using dynamic programming // IEEE TPAMI. 7(2). 1985.
	С. 139–154.
	
	\bibitem{6}
	Veksler O. Stereo correspondence by dynamic programming on a
	tree // Proc. CVPR. Vol. 2. 2005. С. 384–390.
	
	\bibitem{7}
	Naveed I. R., Huijun Di, GuangYou Xu Refine stereo
	correspondence using bayesian network and dynamic programming
	on a color based minimal span tree // ACIVS. 2006. С. 610–619.
	
	\bibitem{8}
	Kolmogorov V., Zabih R. Computing visual correspondence with
	occlusions using graph cuts // ICCV. Vol. 2. 2001. С. 508–515.
	
	\bibitem{9}
	Boykov Y., Veksler O., Zabih R. Fast approximate energy
	minimization via graph cuts // IEEE TPAMI 23(11). 2001.
	С. 1222–1239.
	
	\bibitem{10}
	Sun J., Shum H., Zheng .N Stereo matching using belief
	propagation // In ECCV. 2002. С. 510–524.
	
	%\bibitem{RANSAC}
	%Derpanis K.G. Overview of the RANSAC Algorithm. // 2010 С. 1-2.
	
	\bibitem{RANSAC}
	Fischler M. A. and Bolles R. C. Random sample consensus: a paradigm for model fitting with applications to image analysis and automated cartography. // Commun. ACM 24. 1981. С. 381-395
	
	\bibitem{Sift}
	Lowe D. G. Object recognition from local scale invariant features // Proceedings of the Seventh IEEE International Conference on Computer Vision 2. 1999. С. 1150-1157
	
	\bibitem{Sift2}
	Lowe D. G. Distinctive Image Features from Scale-Invariant Keypoints. // International Journal of Computer Vision. 2004. С. 91-110
\end{thebibliography}



%\appendix   % Тут идут приложения

%\chapter{Картинки}
\label{cha:appendix1}

\blindtext
\begin{figure}
\centering
\caption{Картинка в приложении. Страшная и ужасная.}
\end{figure}

%%% Local Variables: 
%%% mode: latex
%%% TeX-master: "rpz"
%%% End: 


%\chapter{Еще картинки}
\label{cha:appendix2}
\blindtext

\begin{figure}
\centering
\caption{Еще одна картинка, ничем не лучше предыдущей. Но надо же как-то заполнить место.}
\end{figure}

%%% Local Variables: 
%%% mode: latex
%%% TeX-master: "rpz"
%%% End: 


\end{document}

%%% Local Variables:
%%% mode: latex
%%% TeX-master: t
%%% End:
